% ===== §1 Prelude: A Flower on the Path =====
\section{Prelude: A Flower on the Path}
\label{p14:sec:prelude}

\noindent
There is a moment so ordinary that philosophy has almost never paused to look at it. You are walking---not toward anything in particular, or perhaps toward something quite specific, it does not matter---and at the edge of the path there is a flower. It has not been placed there. No one arranged it for you. It is simply there, in the way that contingent things are simply there: arrived from nowhere, held in place by soil and light and the indifference of the season, and likely to be gone before you pass this way again. You notice it. Something in you is touched.

Now: you are alone, and the moment passes into you and is absorbed and becomes part of the texture of an ordinary afternoon. This too is worth philosophical attention, and we will return to it. But it is not the moment that concerns this paper.

The moment that concerns this paper is different. You are not alone. Beside you---or perhaps a little ahead, or half-turned toward something else---is the person you love. And you see the flower, and without thinking, without calculating whether it is worth mentioning, without asking whether she will find it as beautiful as you do, you say: \emph{look}. Or perhaps you say nothing at all, and simply touch her arm and incline your head toward it. She looks. And something happens---not in you, not in her, but \emph{between} you---that is, in some sense that this paper will spend considerable effort trying to specify, \emb{happiness}.

\emc{Why is this different?} The flower has not changed. The path has not changed. The quality of the light, the fragrance if there is one, the precise shade of the petals: all of this is the same whether you are alone or together. And yet the experience is not the same, not in any way that can be accounted for by simple addition---your pleasure plus her pleasure---because what arises when the flower is shared is not a doubled pleasure but something of a different kind altogether. The solitary encounter with the flower is one thing. The shared encounter is not twice that thing. It is another thing.

This observation---so modest as to seem almost trivial---is the seed of the present paper. For it points toward something that neither the existentialist tradition nor the eudaimonist tradition, in their dominant forms, has been equipped to explain. Both traditions have given us rich and sophisticated accounts of what happiness is and how it arises. Both have, in their ways, acknowledged that human beings are not solitary creatures and that relations matter to flourishing. But both have assumed, at the level of their fundamental ontology, that the \emph{subject} of happiness is the individual: that happiness is something that happens \emph{in} a person, is \emph{felt} by a person, is \emph{achieved} by a person---even when that achievement involves and requires others. The flower on the path troubles this assumption. It suggests that what happened in that moment was not, in the first instance, something that happened in you or in her, but something that happened \emph{between} you---in the relational space that your shared attention, your shared presence, your shared life had constituted.

\begin{claim}[The primary philosophical question]
The question this paper addresses is not \emph{what is happiness?} as though happiness were a property of individuals waiting to be correctly analysed. It is: \emph{where does happiness happen?} And the answer it will defend is: in the relational field, as an emergent property of co-evolving dynamical systems, irreducible to the inner state of either participant.
\end{claim}

To defend this answer, the paper must do several things at once. It must engage seriously with the existentialist tradition---with Heidegger's account of thrownness and \emph{Mitsein} \parencite{heidegger1962}, with Sartre's radical contingency \parencite{sartre1943}, with Merleau-Ponty's embodied subjectivity \parencite{merleau-ponty1962}, and with Buber's \emph{I--Thou} \parencite{buber1970}---showing both what these traditions illuminate and where, precisely, they fall short of what the shared flower requires. It must engage equally seriously with the eudaimonist tradition---with Aristotle's account of \emph{eudaimonia} and \emph{phronesis} \parencite{aristotle2000}, with Csikszentmihalyi's psychology of flow \parencite{csikszentmihalyi1990}, with the self-determination theory of \textcite{deci1985}, with Frankl's meaning-centred existentialism \parencite{frankl1959}, and with the contemporary PERMA model \parencite{seligman2011}---showing how each concept is transformed, and in some cases inverted, when the subject of flourishing is relocated from the individual to the relational field.

But the paper must also do something that philosophy alone cannot do. The claim that happiness is an emergent property of co-evolving relational dynamics is not merely a philosophical claim; it is a claim about the structure of dynamical systems, and it can be given formal content---and, importantly, empirical constraint---by drawing on theoretical physics, on computational neuroscience, and on the biology of complex systems. \xref{p14:sec:coevolution} and \xref{p14:sec:dynamics} develop the formal framework, drawing on coupled oscillator theory \parencite{kuramoto1984}, quantum entanglement structure \parencite{einstein1935}, and---in what is perhaps the paper's most original formal contribution---an analogy between the event-driven dynamics of the relational system and the spike-timing dependent plasticity of spiking neural networks \parencite{maass1997,gerstner1996}. \xref{p14:sec:empirical} then translates this framework into a research programme: how would one actually test whether a contingent event descends upon a relational system rather than upon the individuals composing it? What would the EEG and MEG hyperscanning data look like \parencite{hasson2012,liu2023}? What longitudinal designs would detect the co-evolutionary dynamics we are predicting?

The paper proceeds as follows. \xref{p14:sec:existentialism} traces the legacy and limits of existentialism with respect to shared contingency. \xref{p14:sec:eudaimonia} maps the presuppositions of traditional eudaimonism. \xref{p14:sec:turn} executes the relational turn, establishing the GRB framework as the paper's theoretical home. \xref{p14:sec:event} analyses the structure of the contingent event's descent into the relational system, including its double face of joy and suffering. \xref{p14:sec:coevolution} develops the co-evolutionary account of relational dynamics in three registers: philosophical, formal-physical, and neuroscientific. \xref{p14:sec:dynamics} extends this into a full account of developmental dynamics under relational coupling, including the SNN analogy. \xref{p14:sec:empirical} proposes the empirical methodology. \xref{p14:sec:phenomenology} offers a phenomenological description of relational happiness before the formal reconstruction of \xref{p14:sec:reconstruction}. \xref{p14:sec:hermeneutics} re-reads eleven classical happiness concepts under GRB. \xref{p14:sec:culture} situates the account cross-culturally. \xref{p14:sec:praxis} turns to practice. \xref{p14:sec:conclusions} draws the conclusions. The Envoi returns to the flower.

\emc{A word on this paper's place in the series.} \pinkref{Paper~IX} \parencite{paperix} asked how a relation \emph{endures}---how the generative cycle avoids exhaustion and achieves the spiral rather than the circle. \pinkref{Paper~XIII} asked how trust is produced in the absence of external guarantee---how a vow without a sword can generate genuine expectation. The present paper asks a question that is in some sense prior to both: \emph{what is it that makes a relational encounter happy?} Not happy in the thin sense of pleasant, but happy in the eudaimonist sense: constitutive of flourishing, expressive of what it is for a relational being to live well. The answer---that happiness is the signal of generative co-evolution, the felt index of two dynamical systems developing together in response to the contingency that enters their shared field---is, we believe, the missing eudaimonics of the GRB programme: the account of what flourishing \emph{feels like} from the inside of the relational field, and why that feeling is not incidental to flourishing but its very form.
